\documentclass[11pt,a4paper]{article}
\usepackage[T1]{fontenc}
\usepackage[a4paper,margin=2.5cm]{geometry}
\usepackage{url}
\pagestyle{empty}
\setlength{\parindent}{0pt}
\setlength{\parskip}{0.4em}

\begin{document}

\begin{center}
    {\Large\bfseries Individual Contribution Statement}\\[0.3em]
    {\large Project 1 (ODE IVP)}
\end{center}

\textbf{Name:} Tianrui Ji \hfill \textbf{Student ID:} 2361367\\
\textbf{Role:} Algorithms, numerical experiments and visualisation.

\textbf{1. My specific contributions to this project:}

I implemented the matrix ODE model, analytic Jacobian, finite-time exact
solution and numerical diagnostics in \path{Project1/code/model.py}. In
\path{Project1/code/solvers.py}, I implemented Explicit Euler, classical RK4
and Implicit Euler with damped Newton iteration, together with fixed-step
and adaptive integration. I developed the experiments and plotting code in
\path{Project1/code/experiments.py} and \path{Project1/code/run_all.py},
producing the convergence, stability, trajectory, adaptivity, cost and
rank-deficiency results used in the report. For the advanced comparison in
Section 5, I implemented the two-stage Gauss implicit RK4 solver and
four-method comparison in \path{Project1/code_advanced/}, including its four
figures and supporting data. I checked that the numerical claims and figures
used in Sections 3--5 matched the generated results and the theoretical
finite-time solution.

\textbf{2. The hardest conceptual obstacle I encountered:}

The main challenge was keeping the algorithm, theoretical definitions and
reported errors consistent. In particular, a small Newton residual means
that an implicit step equation has been solved closely; it does not imply
an equally small error in the final trajectory. I checked the implemented
residual and Jacobian against the equations, then used the finite-time exact
solution to assess the time-stepping error separately.

\textbf{3. One piece of AI-generated output my team used, and what I changed:}

Appendix A, item AI-2, records AI assistance with numerical-results summaries
and Section 2 wording. I revised an AI-assisted explanation of the $10^{-12}$
Newton tolerance to make clear that it controls the algebraic residual within
each implicit step, not the number of correct digits in the final result.
I checked the revised explanation against the solver and the measured error
relative to the exact solution.

\textbf{4. One thing I still do not fully understand:}

I cannot yet derive a rigorous global nonlinear step-size guarantee from the
frozen Jacobian spectrum. The explicit step-size thresholds in our report
describe local linearised behaviour; extending them to a guarantee for the
full nonlinear trajectory would require further analysis.

\vspace{0.5em}
\textbf{Signature:} \underline{\hspace{4.5cm}} \hfill
\textbf{Date:} \underline{\hspace{2.5cm}}

\end{document}
